\chapter{1927:Heinrich Wieland}

\label{chap:chemistry1927}

The Nobel Prize in Chemistry for the year 1927 was awarded to Heinrich Wieland.

\textbf{Motivation:}

for his investigations of the constitution of the bile acids and related substances

\section{Heinrich Wieland}

\subsection*{Early Life and Education}

Heinrich Otto Wieland was born on 1877-06-04 in Pforzheim, Germany.
He was a German organic chemist awarded the Nobel Prize in Chemistry in 1927 for his investigations of the constitution of the bile acids and related substances.

\subsection*{Career and Major Research}

Wieland studied chemistry at the University of Munich, where he received his doctorate in 1901 under Johannes Thiele.
He taught at the University of Freiburg, where he became professor, and in 1921 he succeeded Richard Willstätter as professor of chemistry at the University of Munich.
His early work concerned nitrogen compounds such as hydrazoic acid, and his studies on the mechanisms of biological oxidation introduced the concept of dehydrogenation.
He also investigated the structures of the alkaloid strychnine, the pigments of butterfly wings (the pterins), and the toxins of the death-cap mushroom.
His most extended research concerned the bile acids, where he established that cholic acid and related compounds share a ring system closely related to cholesterol.

\subsection*{Nobel Prize-winning Work}

The work for which Heinrich Wieland was awarded the Nobel Prize in Chemistry in 1927 was described by the Nobel Committee as: "for his investigations of the constitution of the bile acids and related substances".

Wieland's investigations showed that the bile acids, which occur in bile as salts and are important for the digestion of fats, are built on a common ring skeleton closely related to cholesterol.
This clarified the constitution of an entire family of naturally occurring substances and provided a basis for the later chemistry of the steroid hormones.

\subsection*{Significance and Impact}

By establishing the constitution of the bile acids, Wieland laid the foundations of steroid chemistry.
His concept of biological oxidation as a process of dehydrogenation also became a cornerstone of biochemistry.

\subsection*{Later Career and Honors}

Wieland remained at the University of Munich until his retirement in 1950.
He held honorary doctorates from several universities and was a member of several scientific academies.
He died on 1957-08-05 in Munich, Germany.

\subsection*{Legacy}

The Heinrich Wieland Prize, awarded since 1964 for outstanding research on lipids, is named in his honour.
Wieland's work on the bile acids and on oxidation mechanisms continues to be cited and built upon across the chemical sciences.

\subsection*{Selected Publications}

\begin{itemize}

\item Wieland, H. (1912). "Studien über den Mechanismus der Oxydationsvorgänge." Berichte der deutschen chemischen Gesellschaft.

\item Wieland, H. (1928). "On the Mechanism of Oxidation." Nobel Lecture, Stockholm.

\end{itemize}

\clearpage

\section*{Chapter Summary}

The 1927 prize recognized Heinrich Wieland for his investigations of the constitution of the bile acids and related substances. His work showed that bile acids are derived from a common steroid skeleton and helped establish the chemistry of steroids.
